\section{Introduction}\label{sec:intro}

AlphaFold2 predicts the three-dimensional structure of a protein from its amino-acid sequence~\cite{jumper2021alphafold}.
The public implementation of AlphaFold2 is written in JAX~\cite{alphafoldrepo}. JAX compiles the model with the
accelerated linear algebra (XLA) compiler for central processing units (CPUs), graphics processing units (GPUs), and
tensor processing units (TPUs)~\cite{frostig2018jax}. A TPU slice connects several TPU chips, and JAX treats each chip
as a separate device. A JAX program on a slice does not spread its work over every chip by itself. The developers of
AlphaFold2 state that they do not parallelize the model over several GPUs~\cite{af2issue312} and that the release does
not support running on a TPU~\cite{af2issue611}. The program has to map its work to the chips, for example with
\texttt{jax.pmap}~\cite{jaxpmap}. At list price, a slice costs the same whether one chip works or all eight chips work
(Section~\ref{sec:res-cost}). As a result, a user who runs AlphaFold2 on a slice without mapping the work can pay for
chips that do nothing.

Our goal was to answer two questions for a user of AlphaFold2 on a TPU slice. First, how much of the slice does
AlphaFold2 inference use by default? Second, which way of mapping the work in JAX uses the rest of the slice? For that
way, we also ask how throughput grows with the number of chips.

This study began as a project in a high-performance computing course at Stanford University. TPUs were a main topic.
Afterward, the project grew into the study that we report here.
We ran AlphaFold2 inference on an eight-chip TPU v5e slice of the course cluster in 12 experiments. % src: recomputed: COUNT.tpu_experiments_with_raw_data
We compared four ways of mapping the work. The four ways are the default placement, batching with \texttt{jax.vmap},
multi-query mapping with \texttt{jax.pmap}, and automatic partitioning. The model used random weights, so the study
concerns execution and not the quality of the predicted structures.

Work on AlphaFold2 performance has mostly targeted GPUs.
The authors of AlphaFold2 trained the model on 128 TPU v3 cores and reported inference times on V100 GPUs~\cite{jumper2021alphafold}. % src: jumper2021alphafold ("hence the model uses 128 TPU v3 cores", "using a single model on V100 GPU")
OpenFold reimplements AlphaFold2 in PyTorch~\cite{ahdritz2024openfold}. FastFold runs AlphaFold2 on GPU
clusters~\cite{cheng2024fastfold}. ParaFold splits the pipeline between CPU nodes and GPUs~\cite{zhong2022parafold}.
ParaFold sorts its input sequences by length, so that JAX compiles the model once for many of
them~\cite{zhong2022parafold}.
ColabFold speeds up batch predictions about 90-fold, in part by avoiding recompilation~\cite{colabfold2022}. % src: colabfold2022 ("speeds up batch predictions by ~90-fold by avoiding recompilation and adding an early stop criterion"), paper-v4/SOURCE_NOTES.md section 1
APACE runs AlphaFold2 on 200 NVIDIA A100 GPUs~\cite{apace2024}. % src: apace2024 abstract ("Using up to 300 ensembles, distributed across 200 NVIDIA A100 GPUs"), paper-v4/SOURCE_NOTES.md section 1
ManyFold trains a protein folding model on TPUs~\cite{villegasmorcillo2023manyfold}.
JAXBench includes one AlphaFold2 operator, triangle multiplication, among its 50 TPU workloads~\cite{jaxbench2026}. % src: jaxbench2026 ("JAXBench comprises 50 JAX workloads", "AlphaFold2 triangle multiplication")
A concurrent working paper measures ColabFold inference on one Cloud TPU v5e or v6e chip and AlphaFold2 model
inference on one v6e chip~\cite{forsyth2026timetoscience}.
On November 23, 2026, we searched OpenAlex, DBLP, arXiv, IEEE Xplore, Crossref, Europe PMC, the Google Cloud blog, and the GitHub issues of AlphaFold2 and ColabFold. % src: paper-v4/LITERATURE_SEARCH.md ("23 November 2026")
Apart from the earlier versions of this paper\ifappendix{} (Appendix~\ref{app:changes})\fi, we found no published
measurement of AlphaFold2 inference across a multi-chip TPU slice.

In this paper, a query is one input sequence that the model predicts once. Besides the experiments on the slice, we ran
one baseline configuration (Section~\ref{sec:workload}) on one TPU chip, on an NVIDIA T4 GPU, and on a Google Colab
CPU runtime.

By default, AlphaFold2 ran on one chip of the slice, and the other seven chips held no data. Batching up to eight
queries on one chip with \texttt{jax.vmap} did not raise throughput. With \texttt{jax.pmap}, the eight chips ran eight
independent queries in parallel.
On a grid of chip counts and sequence lengths, the full slice gave 6.53 to about 7.9 times the throughput of one chip. % src: results/sweep/scaling_law_data.json -> grid; recomputed: V2.speedup_8over1_min, V2.speedup_8over1_max
Without sharding annotations, JAX's automatic partitioning did not shorten the steady-state call for one query. We also
built an external ensemble, which runs differently featurized copies of one query outside the model. The external
ensemble ran on all eight chips. A compilation cache shared by two processes shortened parameter initialization and
the first call of the second process. Section~\ref{sec:timing} defines these steps.

We report the comparison with the T4 and with the CPU runtime as context and not as a result. The three platforms ran
in separate sessions and with different copies of the benchmark script. On the T4 and on the CPU runtime, the times
changed between sessions (Section~\ref{sec:res-platforms}). The results on the slice come from one platform and from
one period.

% \keepwithlist (macros.tex) keeps the next sentence with the first item of the list, so that it never ends a column.
\keepwithlist
This paper makes three contributions.
\begin{itemize}
\item We measure how AlphaFold2 inference uses an eight-chip TPU v5e slice under the four ways of mapping the work
described above (Sections~\ref{sec:res-default} to~\ref{sec:res-auto}).
\item We give the throughput on a grid of chip counts and sequence lengths when each chip runs one query through
\texttt{jax.pmap} (Section~\ref{sec:res-pmap}). We also give the cost of the queries at list price
(Section~\ref{sec:res-cost}).
\item We list the script, the input, the timed region, and the number of runs of every experiment. We also state
which records were not kept (Section~\ref{sec:method}).
\end{itemize}

% Version 4 of arXiv:2609.34818 (December 2026).
This paper is version 4 of arXiv:2609.34818.
\ifappendix Appendix~\ref{app:changes} lists the claims that changed from the earlier versions.\fi
